\chapter{From a physical story to a mathematical claim}\label{ch:contract}
A store has more water than its declared reference and a clinic has less.
Moving water from one to the other seems helpful. Mathematics begins by
asking what that sentence has left unspecified. How much water? Which
reference? Which route? What else changes? Does the action remain helpful
when it is large? Who is permitted to perform it? The value of a formal
model is not that it makes these questions disappear. It gives each question
an object precise enough to inspect.

\section{The declared world}
Let there be a finite number $n$ of accounting cells, with state
\begin{equation}
 x=(x_1,\ldots,x_n)^T\in\mathcal X_M,
 \qquad \mathcal X_M=\{x\in\mathbb R^n:x_i\geq0,\ \sum_i x_i=M\}.
\end{equation}
Every $x_i$ measures the same homogeneous resource. The superscript $T$
turns a written row into a column. The set $\mathcal X_M$ is the nonnegative
fixed-mass simplex: all ways of distributing the same total $M$ among the
cells. For two cells it is a line segment, not the whole plane. For three it
is a triangle embedded in three-dimensional space.

The programme declares a reference $x^*\in\mathcal X_M$ and scales
$\sigma_i>0$. Membership of the reference in the same simplex matters.
If two cells contain twenty units in total, references requiring twelve
units in each cannot be reached by closed internal transfer. That mismatch
is not an actor failure. It is a different modelling problem.

Compatibility is necessary, not sufficient, for reachability. A graph with
no route to the clinic cannot deliver water there. A menu containing only
four-unit transfers may be unable to remove a one-unit discrepancy exactly.
The geometry says which states are allowed by total mass; the action system
says which transitions are available.

\section{Definition, proposition and evidence}
We define a potential and derive its consequences. A definition is a choice
about the model, not a discovery that nature uses that formula. A proposition
may then be true for every state in the declared domain. A worked example
can help a reader follow the proof but cannot replace it. An experiment
would examine a specified transition process under registered conditions.

The programme document \cite{programme} makes this distinction explicit.
The active model is prospective. Its exact quadratic identities can be
proved now. Whether a random-affordable actor process produces useful
recovery, distribution or long-run behavior is a different question.

\begin{lawbox}{Standing assumptions for active examples}
One homogeneous conserved scalar; a finite number of cells; nonnegative
stocks; fixed compatible references; fixed positive scales; lossless internal
actions; additive declared increments. Unless a section explicitly says
otherwise, separate process burden is zero. No restoring dynamics or
optimizing actor is hidden in these assumptions.
\end{lawbox}

\section{What the state leaves outside}
A simple water stock model does not automatically represent water quality,
pumping energy, soil damage or a patient's health. These omissions limit the
claim. They cannot be repaired by saying the scalar potential is holistic.
If quality is consequential, it needs a represented coordinate or an
explicitly scoped limitation. If pumping consumes electricity, a broader
physical account must link that transformation to an energy carrier.

Nor does the model identify a reference with a moral optimum. A reference
may later be informed by functional evidence and public choices, but the
proof that follows only knows its declared number. Keeping this boundary
visible lets mathematical work proceed without pretending that calibration
or governance has already been solved.

\section{A useful proof-reading habit}
Before accepting a theorem, write its domain in words. When we later prove
that $V$ is bounded, the reason will be finite mass and fixed scales, not an
EBU controller. When we prove that a sequence telescopes, the reason will be
consecutive endpoint cancellation, not benevolent actors. When we prove that
path contributions sum to the group value, the reason will be linearity and
the fundamental theorem of calculus, not a fair ownership rule.

This habit prevents a correct sentence from acquiring an incorrect sequel.
``The accounts close'' does not imply ``the world recovers.'' ``The action
has positive value'' does not imply ``the action is physically possible.''
Each sequel needs its own premises.

\section*{Exercises and discussion}
For a three-cell world with mass twelve, decide whether $(4,4,4)$,
$(5,5,5)$ and $(12,0,0)$ are compatible references. The first and third are;
the second is not. Compatibility of the third still says nothing about
whether the graph can move all stock to cell one.

Explain why the phrase ``Gaussian equilibrium'' is potentially misleading.
The Gaussian construction defines reference geometry. To call the reference
an equilibrium of a process, one must specify that process and show that it
leaves the reference unchanged. To call it attracting requires more still.

\section*{Chapter recap}
A model needs a domain, a state, a reference, a transformation and a claim.
The rest of this volume makes those objects explicit. Its mathematical
results are conditional statements about the declared objects, not a
shortcut around physical measurement or scientific testing.
